\documentclass[11pt,a4paper]{article}
\usepackage[margin=22mm]{geometry}
\usepackage{amsmath,amssymb}
\usepackage{fontspec}
\usepackage{polyglossia}
\setdefaultlanguage{russian}
\setmainfont{DejaVu Serif}
\usepackage[unicode,hidelinks]{hyperref}

\allowdisplaybreaks[1]
\raggedright
\setlength{\parskip}{0.35em}
\newcommand{\dd}{\mathop{}\!\mathrm{d}}
\newcommand{\avg}[1]{\langle #1\rangle}
\newcommand{\ket}[1]{\lvert #1\rangle}
\newcommand{\bra}[1]{\langle #1\rvert}
\newcommand{\braket}[2]{\langle #1\mid #2\rangle}
\DeclareMathOperator{\Var}{Var}

\newcommand{\solutiontitle}[4]{%
  \hypersetup{pdftitle={Квантовая механика 1. Решения ДЗ №#1}}%
  \begin{center}
    {\small Квантовая механика 1 \enspace\textperiodcentered\enspace Семестр 1\par}
    \vspace{1ex}
    {\LARGE\bfseries Решения домашнего задания №#1\par}
    \vspace{0.5ex}
    {\large #2\par}
  \end{center}
  \noindent{\small Условия приведены в файле «#3», #4.}\par\medskip
}

\begin{document}
\solutiontitle{3}{Вероятности, чётность и средние значения}{ДЗ2026(1)\_03.pdf}{с.~1--2}

\section{Спектральное вычисление моментов}

Берём \(\psi=(1,-i,i)^{\mathsf T}/\sqrt3\). При
\(v=(0,1,i)^{\mathsf T}/\sqrt2\) заданная матрица удовлетворяет
\(Ke_1=\sqrt2 v,\ Kv=\sqrt2 e_1\), а
\(K(0,i,1)^{\mathsf T}=0\). Нормированные собственные векторы
\(u_\pm=(e_1\pm v)/\sqrt2,\ u_0=(0,i,1)^{\mathsf T}/\sqrt2\)
отвечают \(\lambda_\pm=\pm\sqrt2,\ \lambda_0=0\).
Соответственно
\[
 p_\pm=\left|\frac1{\sqrt6}\pm\frac{1-i}{2\sqrt3}\right|^2
       =\frac13\pm\frac1{3\sqrt2},
 \qquad p_0=\left|\frac{-1+i}{\sqrt6}\right|^2=\frac13.
\]
Все три исхода возможны, поскольку все три вероятности положительны.
По формулам моментов дискретного распределения
\[
 \avg{K}=\sqrt2(p_+-p_-)=\frac23,\qquad
 \avg{K^2}=2(p_++p_-)=\frac43,\qquad \Var(K)=\frac89 .
\]

\section{Измерение матрицы Паули}

Матрица \(\sigma_x\) переставляет две компоненты. Её нормированные
собственные векторы \(v_\pm=(1,\pm1)^{\mathsf T}/\sqrt2\)
имеют собственные значения \(\pm1\).
Для \(\psi=(1,-2i)^{\mathsf T}/\sqrt5\) проекции равны
\(\braket{v_\pm}{\psi}=(1\mp2i)/\sqrt{10}\), значит
\(p_\pm=(1+4)/10=1/2\).

\section{Одномерное состояние с гауссовой огибающей}

Положим \(y=x/\sigma,\ \sigma>0\). Плотность для \(y\)
равна \(w(y)=|C|^2(1+y)^2e^{-y^2}/\sqrt\pi\).
Используем гауссовы моменты
\[
 G_j=\frac1{\sqrt\pi}\int_{\mathbb R}y^j e^{-y^2}\,\dd y,\qquad
 G_0=1,\quad G_2=\tfrac12,\quad G_4=\tfrac34,\quad G_{2j+1}=0.
\]
Чётные моменты получаются по частям из
\(G_{2j}=(2j-1)G_{2j-2}/2\), нечётные обнуляются по симметрии.

\textbf{(a)} Норма равна \(|C|^2(G_0+2G_1+G_2)=3|C|^2/2\).
Выбираем \(C=\sqrt{2/3}\).

\textbf{(b)} Имеем \(\avg{y}=(2/3)(2G_2)=2/3\) и
\(\avg{y^2}=(2/3)(G_2+G_4)=5/6\). Значит
\[
 \avg{x}=\frac{2\sigma}{3},\qquad
 \Var(x)=\sigma^2(5/6-4/9)=\frac{7\sigma^2}{18}.
\]

\textbf{(c)} Производная \(w'(y)\) пропорциональна
\(e^{-y^2}(1+y)[1-y(1+y)]\). Для \(y>0\)
\(w(y)>w(-y)\), а \(w'(0)>0\). Поэтому максимум следует искать
на положительной полуоси, где производная меняет знак с плюса на
минус в единственном корне \(1-y-y^2=0\).
Искомая координата \(x_{\rm mp}=\sigma(\sqrt5-1)/2\).

\textbf{(d)} В \(\Psi(x)^*\Psi(-x)\) полиномиальный множитель
пропорционален \(1-y^2\). Получаем
\(\avg{P}=(2/3)(G_0-G_2)=1/3\). Так как \(P^2=I\),
\(\Var(P)=1-(1/3)^2=8/9\).

\textbf{(e)} Чётная часть плотности даёт на отрицательной
полуоси ровно \(1/2\). Нечётная часть даёт
\((4/3\sqrt\pi)\int_{-\infty}^0 y e^{-y^2}\dd y=-2/(3\sqrt\pi)\).
Таким образом \(\Pr(x\le0)=1/2-2/(3\sqrt\pi)\).

\section{Сферически-симметричная экспонента}

Для \(a>0\) элемент объёма равен \(r^2\dd r\,\dd\Omega\).
По формуле \(\int_0^\infty r^n e^{-br}\dd r=n!/b^{n+1}\),
получаемой \(n\) интегрированиями по частям,
\(4\pi C^2\int_0^\infty r^2 e^{-2r/a}\dd r=\pi a^3 C^2=1\).
Берём \(C=(\pi a^3)^{-1/2}\). Радиальная плотность относительно
\(\dd r\) есть \(w_r(r)=4r^2e^{-2r/a}/a^3\).

\textbf{(a)} По той же формуле
\[
 \avg{r^n}=\frac{(n+2)!}{2}\left(\frac a2\right)^n,\qquad
 \avg{r}=\frac{3a}{2},\qquad \avg{r^2}=3a^2,\qquad
 \Var(r)=\frac{3a^2}{4}.
\]

\textbf{(b, c)} Плотность инвариантна при \(\mathbf r\mapsto-\mathbf r\),
и каждая компонента \(\mathbf r\) меняет знак, так что
\(\avg{\mathbf r}=\mathbf 0\). Сама функция тоже чётна, значит
\(P\Psi=\Psi\) и \(\avg{P}=1\).

\textbf{(d)} Для слоя \([a/2,a/2+\Delta]\) точная вероятность
равна \(\int_{a/2}^{a/2+\Delta} w_r(r)\dd r\). При \(u=2r/a\)
первообразная равна \(-e^{-u}(1+u+u^2/2)\).
Так как \(w_r(a/2)=1/(ea)\), для \(\Delta\ll a\) получаем
\(\Pr(\text{слой})=\Delta/(ea)+O((\Delta/a)^2)\).
Если слой центрирован на \(a/2\), его вероятность имеет тот же
ведущий член для полной толщины \(\Delta\).

\textbf{(e)} Отражение \(z\mapsto-z\) меняет полупространства
местами и сохраняет плотность. Плоскость \(z=0\) имеет нулевую
вероятность, поэтому \(\Pr(z\le0)=1/2\).

\section{Ядра и теорема Геллмана --- Фейнмана (5*)}

\textbf{(a)} Дифференцируя по \(x\) тождество
\(\int\delta(x-x') f(x')\dd x'=f(x)\), находим ядро
\(-i\partial_x\delta(x-x')\). Для чётности свойство
\(\int\delta(x+x')f(x')\dd x'=f(-x)\) даёт ядро \(\delta(x+x')\).

\textbf{(b)} Пусть нормированная ветвь \(\ket{n(\lambda)}\) и её
собственное значение дифференцируемы, а вектор и его производная
лежат в области, позволяющей дифференцировать собственное уравнение.
Опуская аргумент \(\lambda\), имеем
\[
 F'\ket n+F\ket{n'}=f'_n\ket n+f_n\ket{n'}.
\]
Умножаем слева на \(\bra n\). Из самосопряжённости и вещественности
\(f_n\) следует \(\bra n F=f_n\bra n\), так что оба слагаемых
\(f_n\braket{n}{n'}\) сокращаются. Поскольку \(\braket{n}{n}=1\),
\[
 \frac{\dd f_n}{\dd\lambda}
 =\bra{n(\lambda)}\frac{\partial F}{\partial\lambda}\ket{n(\lambda)}.
\]
При вырождении доказательство работает для любой выбранной
дифференцируемой собственной ветви.

\end{document}
